\section{Gluon quasi-PDF operators}\label{sec:operators}

Now we are ready to construct an appropriate gluon quasi-PDF operator.
To this end, we need to identify one of the indices in $J_1^{\mu\nu}$ with $z$ or $t$ and let the other run either over all Lorentz components or over transverse components only. It is worthwhile to point out at this stage that the operator $J_3^{\mu\nu}$ only yields contact terms when integrating out the ``heavy quark" field. This can be seen from the fact that the EOM operator acting on the ``heavy quark" propagator yields a $\delta$-function. The contact terms do not vanish only when $z_2=z_1$, indicating that an extra renormalization is required when two LCOs shrink to one spacetime point. For the renormalization of the non-local gluon quasi-PDF operator, $J_3^{\mu\nu}$ is irrelevant and can be ignored.

From the discussions above, we identify the following building blocks, $J_{1,R}^{zi}$, $J_{1,R}^{ti}$, $J_{1,R}^{z\mu}$, for the gluon quasi-PDF operator. A multiplicatively renormalizable gluon unpolarized quasi-PDF operator can be constructed from
\begin{align}
 {\cal O}^{1}_R(z_2,z_1) &\equiv  J_{1,R}^{ti}(z_2) \overline J_{1,R}^{ti}(z_1), \non\\
 {\cal O}^{2}_R(z_2,z_1) &\equiv  J_{1,R}^{zi}(z_2)\overline J_{1,R}^{zi}(z_1), \nonumber\\
 {\cal O}^{3}_R(z_2, z_1) &\equiv    J_{1,R}^{ti}(z_2) \overline J_{1,R}^{zi}(z_1),\non\\
 {\cal O}^{4}_R(z_2, z_1)  &\equiv  J_{1,R}^{z\mu}(z_2) \overline J_{1,R,\mu}^{z}(z_1).
\end{align}
After integrating out the auxiliary ``heavy quark" field~\cite{Ji:2017oey}, these operators renormalize multiplicatively as
\begin{align}
{O}^{1}_R(z_2, z_1) &=Z_{11}^2 e^{\overline{\delta { m}}\Delta z}F^{ti}(z_2) L(z_2,z_1) F^{ti}(z_1),\nonumber\\
 {  O}^{2}_R(z_2,z_1) &=Z_{22}^2 e^{\overline{\delta { m}}\Delta z} F^{zi}(z_2) L(z_2,z_1) F^{zi}(z_1),\nonumber\\
 {O}^{3}_R(z_2,z_1) &= Z_{11}Z_{22} e^{\overline{\delta { m}}\Delta z} F^{ti}(z_2) L(z_2,z_1) F^{zi}(z_1),\nonumber\\
O^{4}_R(z_2,z_1)
& = Z_{22}^2 e^{\overline{\delta { m}}\Delta z} F^{z\mu}(z_2) L(z_2,z_1) F^{z}_{\;\;\mu}(z_1),
 \end{align}
with $\Delta z=|z_2-z_1|$.

In the large momentum limit,
the operators $O^{i}_R(i=1,2,3,4)$ differ from each other only by power corrections. They therefore belong to the same universality class defining the gluon quasi-PDF.
Given the above four operators, one can use any combination of them to study the gluon quasi-PDF, however such combinations are usually not multiplicatively renormalizable. A notable example is
\begin{align}
{  O}^{5}_R(z_2,z_1)
 %\equiv   J_{1,R}^{t\mu}(z_2) \overline J^{t}_{1,R,\mu}(z_1)\non\\
 =    -  {  O}^{1}_R(z_2,z_1)-  {  O}^2_R(z_2, z_1)-{  O}^4_R(z_2, z_1),
\end{align}
which (minus the trace term) has been used in the recent simulation~\cite{Fan:2018dxu}.
However, since the renormalizations for ${  O}^{1}_R(z_2,z_1)$ and  ${  O}_R^{2,4}(z_2, z_1)$ are different, ${  O}_R^{5}(z_2,z_1)$ is not multiplicatively renormalizable.

In the above discussion, we considered the gluon quasi-PDF only. If we insert the gluon quasi-PDF operator into a quark state, it gives rise to UV finite contributions, provided that all subdivergences have been renormalized. The reason is that, the quasi-PDF is defined at spacelike separations, therefore there is no non-local UV divergence apart from the exponential mass renormalization. It indicates that the gluon quasi-PDF does not mix with the quark ones under renormalization in the above sense, while the mixing occurs at the factorization stage. This has been confirmed by the one-loop calculations in Ref.~\cite{Wang:2017qyg}, where the quark matrix element of the gluon quasi-PDF operator is UV finite, but contains $\ln P_z^2/p^2$ ($p^2$ is an IR regulator) terms associated with the corresponding splitting kernel. This will be matched to the quark PDF to ensure the cancellation of IR divergences between the quasi-PDF and the PDF.
